\documentclass[tikz,border=6pt]{standalone}
\usepackage{ctex}
\usepackage{amsmath}
\usepackage{amssymb}
\usepackage{pgfplots}
\pgfplotsset{compat=1.18}
\usetikzlibrary{calc,angles,quotes,arrows.meta,positioning,decorations.markings}
\begin{document}
\begin{tikzpicture}

% ===== 左图：单位圆上的两个终边解 =====
\begin{scope}[scale=2.0]
  % 坐标轴
  \draw[->,gray!70] (-1.35,0) -- (1.45,0) node[right,black]{$x$};
  \draw[->,gray!70] (0,-1.35) -- (0,1.45) node[above,black]{$y$};
  % 单位圆
  \draw[thick] (0,0) circle (1);
  % y = 1/2 水平虚线
  \draw[blue!60!black,dashed] (-1.25,0.5) -- (1.25,0.5);
  \node[blue!60!black,font=\small] at (1.06,0.66) {$y=\tfrac12$};
  % 两个解的终边
  \coordinate (A) at (30:1);   % pi/6
  \coordinate (B) at (150:1);  % 5pi/6
  \draw[->,red!75!black,very thick] (0,0) -- (A);
  \draw[->,red!75!black,very thick] (0,0) -- (B);
  % 标记交点
  \fill[red!75!black] (A) circle (0.035);
  \fill[red!75!black] (B) circle (0.035);
  % 纵坐标投影
  \draw[red!60,dotted] (A) -- (0,0.5);
  \draw[red!60,dotted] (B) -- (0,0.5);
  % 角度弧
  \draw[red!75!black] (0.42,0) arc (0:30:0.42);
  \node[red!75!black,font=\small] at (0.62,0.16) {$\tfrac{\pi}{6}$};
  \draw[red!75!black] (0.30,0) arc (0:150:0.30);
  \node[red!75!black,font=\small] at (-0.40,0.30) {$\tfrac{5\pi}{6}$};
  % 终边解标签（移到圆外侧，避开 y=1/2 标签与坐标轴顶部）
  \node[red!75!black,font=\small] at (1.18,0.95) {$x=\tfrac{\pi}{6}$};
  \node[red!75!black,font=\small] at (-1.18,0.95) {$x=\tfrac{5\pi}{6}$};
  % 标题
  \node[font=\bfseries] at (0,-1.62) {单位圆：两个终边解};
\end{scope}

% ===== 右图：sin 曲线与 y=1/2 的无穷多交点 =====
\begin{scope}[xshift=8.0cm]
  % 标题（置于绘图区上方，避开曲线峰值）
  \node[font=\bfseries, anchor=south] at (5.5cm,4.0cm)
    {正弦曲线与 $y=\tfrac12$ 的无穷多交点};
  \begin{axis}[
      width=12cm, height=6.2cm,
      title style={yshift=4pt},
      axis lines=middle,
      xlabel={$x$}, ylabel={$y$},
      xtick={-3.14159,0,3.14159,6.28319,9.42478},
      xticklabels={$-\pi$,$0$,$\pi$,$2\pi$,$3\pi$},
      ytick={-1,0.5,1}, yticklabels={$-1$,$\tfrac12$,$1$},
      xmin=-3.8, xmax=10.4, ymin=-1.55, ymax=1.5,
      enlargelimits=false,
      tick label style={font=\small},
      label style={font=\small},
      legend style={at={(0.5,-0.14)}, anchor=north, legend columns=2,
                    font=\small, draw=none, fill=none, column sep=1.2em},
    ]
    % sin 曲线
    \addplot[domain=-3.8:10.4, samples=300, thick, black] {sin(deg(x))};
    \addlegendentry{$y=\sin x$}
    % y = 1/2 水平线
    \addplot[domain=-3.8:10.4, samples=2, blue!60!black, dashed, thick] {0.5};
    \addlegendentry{$y=\tfrac12$}
    % 交点：pi/6 + 2k*pi 与 5pi/6 + 2k*pi
    \foreach \k in {-1,0,1} {
      \edef\xa{0.523599 + 6.283185*\k} % pi/6 + 2k pi
      \edef\xb{2.617994 + 6.283185*\k} % 5pi/6 + 2k pi
      \edef\temp{\noexpand\addplot[only marks, mark=*, red!75!black, mark size=1.6pt] coordinates {(\xa,0.5)};}
      \temp
      \edef\temp{\noexpand\addplot[only marks, mark=*, red!75!black, mark size=1.6pt] coordinates {(\xb,0.5)};}
      \temp
    }
  \end{axis}
  % 标注两族解（半透明框，置于图例下方空白处，避开峰值与曲线）
  \node[font=\small, fill=white, fill opacity=0.78, text opacity=1, inner sep=2.5pt,
        align=center, anchor=north] at (5.5cm,-2.7cm)
    {通解：$x=\tfrac{\pi}{6}+2k\pi$ \ 或 \ $x=\tfrac{5\pi}{6}+2k\pi,\quad k\in\mathbb{Z}$};
\end{scope}

\end{tikzpicture}
\end{document}
